% open-logic.tex
% compiles to produce complete text using open-logic-dev style

\documentclass[../include/open-logic-part]{subfiles}

\begin{document}

\clearpage

\begin{editorial}
यह फ़ाइल ओपन लॉजिक प्रोजेक्ट में सम्मिलित सारी सामग्री लोड करती है। यदि इस
तरह की संपादकीय टिप्पणियाँ दिखाई दें, तो वे बताती हैं कि इस सामग्री को कैसे
प्रस्तुत किया जाएगा, इसका विचार किए बिना फ़ाइल संकलित की गई थी। यदि आप इसे
पढ़ सकते हैं, तो संभवतः इस पीडीएफ़ से पढ़ाना या पढ़ना \emph{उचित नहीं} है।

ओपन लॉजिक प्रोजेक्ट ऐसी अनेक व्यवस्थाएँ देता है जिनसे अध्यापन या स्वाध्याय
के लिए अधिक उपयुक्त पाठ बनाया जा सकता है। उदाहरणतः पाठ डिफ़ॉल्ट रूप से सभी
तार्किक संकारकों को आदिम मानेगा और प्रमाणों में सभी संकारकों की सभी स्थितियाँ
पूरा करेगा। किंतु इनमें से कुछ स्थितियों को अभ्यास के रूप में छोड़ना कहीं
बेहतर है। ओपन लॉजिक प्रोजेक्ट पर काम भी जारी है। सहयोग और सुधार को प्रोत्साहन
देने के लिए सामग्री तब भी सम्मिलित की जाती है जब वह केवल मसौदा हो, उसमें
अभ्यास न हों, इत्यादि। किसी पाठ्यक्रम के लिए बनी पीडीएफ़ इन खंडों को बाहर
रखेगी।

अध्यापन और अध्ययन के लिए अधिक उपयुक्त पीडीएफ़ पाने हेतु
\href{http://builds.openlogicproject.org/}{ओएलपी वेबसाइट पर उपलब्ध नमूना
पाठ्यक्रम} देखें। अपना पाठ बनाने के लिए आप
\href{https://github.com/OpenLogicProject/OpenLogic/tree/master/courses/sample}{नमूना
ड्राइवर फ़ाइल} से आरंभ कर सकते हैं, या अधिक सुसज्जित और उन्नत उदाहरणों के
लिए व्युत्पन्न पाठ्यपुस्तकों के स्रोत देख सकते हैं।
\end{editorial}

\olimport[sets-functions-relations]{sets-functions-relations-complete}

\olimport[propositional-logic]{propositional-logic}

\olimport[first-order-logic]{first-order-logic}

\olimport[model-theory]{model-theory}

\olimport[computability]{computability}

\olimport[turing-machines]{turing-machines}

\olimport[incompleteness]{incompleteness}

\olimport[second-order-logic]{second-order-logic}

\olimport[lambda-calculus]{lambda-calculus}

\olimport[many-valued-logic]{many-valued-logic}

\olimport[normal-modal-logic]{normal-modal-logic}

\olimport[applied-modal-logic]{applied-modal-logic}

\olimport[intuitionistic-logic]{intuitionistic-logic}

\olimport[counterfactuals]{counterfactuals}

\olimport[set-theory]{set-theory}

\olimport[methods]{methods}

\olimport[history]{history}

\olimport[reference]{reference}

\end{document}

